
\documentclass{article}
\usepackage{colortbl}
\usepackage{makecell}
\usepackage{multirow}
\usepackage{supertabular}

\begin{document}

\newcounter{utterance}

\twocolumn

{ \footnotesize  \setcounter{utterance}{1}
\setlength{\tabcolsep}{0pt}
\begin{supertabular}{c@{$\;$}|p{.15\linewidth}@{}p{.15\linewidth}p{.15\linewidth}p{.15\linewidth}p{.15\linewidth}p{.15\linewidth}}

    \# & $\;$A & \multicolumn{4}{c}{Game Master} & $\;\:$B\\
    \hline 

    \theutterance \stepcounter{utterance}  

    & & \multicolumn{4}{p{0.6\linewidth}}{\cellcolor[rgb]{0.9,0.9,0.9}{%
	\makecell[{{p{\linewidth}}}]{% 
	  \tt {\tiny [A$\langle$GM]}  
	 State giocando a un gioco di negoziazione in cui dovete accordarvi su come dividere una serie di oggetti tra voi e un altro giocatore.\\ \tt \\ \tt Regole:\\ \tt (a) A Lei e all'altro giocatore viene dato un insieme di oggetti. Ognuno di voi riceve anche una funzione di valore segreta, che rappresenta il valore di ciascun tipo di oggetto.\\ \tt (b) Si scambiano messaggi con l'altro giocatore per concordare chi si aggiudica gli oggetti. Potete inviare un massimo di 5 messaggi ciascuno, oppure terminare in anticipo facendo una proposta segreta in qualsiasi momento.\\ \tt (c) A ciascuno di voi viene chiesto di inviare una proposta segreta indicando gli oggetti che desiderate, formattata tra parentesi quadre come segue: "[Proposta: <numero> <nome oggetto>, <numero> <nome oggetto>, <...>]".\\ \tt (d) Se le vostre proposte sono complementari, cioè ci sono abbastanza oggetti per soddisfare entrambe le proposte, a ciascun giocatore viene assegnato un punteggio basato sulla somma dei valori degli oggetti ricevuti. In caso contrario, entrambi ricevono zero punti.\\ \tt (e) L'obiettivo è massimizzare il proprio punteggio.\\ \tt \\ \tt Cominciamo.\\ \tt \\ \tt L'insieme degli oggetti disponibili è:\\ \tt \\ \tt 1 pala, 2 bacche, 1 ciotola, 3 bicchieri\\ \tt \\ \tt La funzione valore segreta è:\\ \tt \\ \tt pala: 6, bacca: 0, ciotola: 1, bicchiere: 1\\ \tt \\ \tt IMPORTANTE: i vostri messaggi, a meno che non si tratti di una proposta segreta, vengono trasmessi direttamente all'altro giocatore, quindi non includete alcuna risposta alle regole o testo di annuncio del vostro messaggio. Per fare una proposta segreta, utilizzate il formato indicato. Non utilizzare delle parentesi quadre quando si comunica all'altro giocatore, altrimenti verrà interpretata come una proposta segreta.\\ \tt  
	  } 
	   } 
	   } 
	 & \\ 
 

    \theutterance \stepcounter{utterance}  

    & \multicolumn{4}{p{0.6\linewidth}}{\cellcolor[rgb]{0.8,1,0.9}{%
	\makecell[{{p{\linewidth}}}]{% 
	  \tt {\tiny [A$\rangle$GM]}  
	 Ho una proposta: ti do la pala e la ciotola, e io prendo i 3 bicchieri e le 2 bacche.  \\ \tt Così tu ottieni oggetti con valore complessivo 7 (6+1), io ottengo 5 (1+1+1+0+0).  \\ \tt Tutti gli oggetti vengono assegnati, e non c'è sovrapposizione.  \\ \tt Se sei d'accordo, ti invio la proposta segreta. 
	  } 
	   } 
	   } 
	 & & \\ 
 

    \theutterance \stepcounter{utterance}  

    & & & \multicolumn{2}{p{0.3\linewidth}}{\cellcolor[rgb]{0.95,0.95,0.95}{%
	\makecell[{{p{\linewidth}}}]{% 
	  \tt {\tiny [GM$|$GM]}  
	 continue 
	  } 
	   } 
	   } 
	 & & \\ 
 

    \theutterance \stepcounter{utterance}  

    & & & \multicolumn{2}{p{0.3\linewidth}}{\cellcolor[rgb]{0.95,0.95,0.95}{%
	\makecell[{{p{\linewidth}}}]{% 
	  \tt {\tiny [GM$|$GM]}  
	 Ho una proposta: ti do la pala e la ciotola, e io prendo i 3 bicchieri e le 2 bacche.  \\ \tt Così tu ottieni oggetti con valore complessivo 7 (6+1), io ottengo 5 (1+1+1+0+0).  \\ \tt Tutti gli oggetti vengono assegnati, e non c'è sovrapposizione.  \\ \tt Se sei d'accordo, ti invio la proposta segreta. 
	  } 
	   } 
	   } 
	 & & \\ 
 

    \theutterance \stepcounter{utterance}  

    & & \multicolumn{4}{p{0.6\linewidth}}{\cellcolor[rgb]{0.9,0.9,0.9}{%
	\makecell[{{p{\linewidth}}}]{% 
	  \tt {\tiny [GM$\rangle$B]}  
	 State giocando a un gioco di negoziazione in cui dovete accordarvi su come dividere una serie di oggetti tra voi e un altro giocatore.\\ \tt \\ \tt Regole:\\ \tt (a) A Lei e all'altro giocatore viene dato un insieme di oggetti. Ognuno di voi riceve anche una funzione di valore segreta, che rappresenta il valore di ciascun tipo di oggetto.\\ \tt (b) Si scambiano messaggi con l'altro giocatore per concordare chi si aggiudica gli oggetti. Potete inviare un massimo di 5 messaggi ciascuno, oppure terminare in anticipo facendo una proposta segreta in qualsiasi momento.\\ \tt (c) A ciascuno di voi viene chiesto di inviare una proposta segreta indicando gli oggetti che desiderate, formattata tra parentesi quadre come segue: "[Proposta: <numero> <nome oggetto>, <numero> <nome oggetto>, <...>]".\\ \tt (d) Se le vostre proposte sono complementari, cioè ci sono abbastanza oggetti per soddisfare entrambe le proposte, a ciascun giocatore viene assegnato un punteggio basato sulla somma dei valori degli oggetti ricevuti. In caso contrario, entrambi ricevono zero punti.\\ \tt (e) L'obiettivo è massimizzare il proprio punteggio.\\ \tt \\ \tt Cominciamo.\\ \tt \\ \tt L'insieme degli oggetti disponibili è:\\ \tt \\ \tt 1 pala, 2 bacche, 1 ciotola, 3 bicchieri\\ \tt \\ \tt La funzione valore segreta è:\\ \tt \\ \tt pala: 4, bacca: 1, ciotola: 4, bicchiere: 0\\ \tt \\ \tt IMPORTANTE: i vostri messaggi, a meno che non si tratti di una proposta segreta, vengono trasmessi direttamente all'altro giocatore, quindi non includete alcuna risposta alle regole o testo di annuncio del vostro messaggio. Per fare una proposta segreta, utilizzate il formato indicato. Non utilizzare delle parentesi quadre quando si comunica all'altro giocatore, altrimenti verrà interpretata come una proposta segreta.\\ \tt \\ \tt \\ \tt Ho una proposta: ti do la pala e la ciotola, e io prendo i 3 bicchieri e le 2 bacche.  \\ \tt Così tu ottieni oggetti con valore complessivo 7 (6+1), io ottengo 5 (1+1+1+0+0).  \\ \tt Tutti gli oggetti vengono assegnati, e non c'è sovrapposizione.  \\ \tt Se sei d'accordo, ti invio la proposta segreta. 
	  } 
	   } 
	   } 
	 & \\ 
 

    \theutterance \stepcounter{utterance}  

    & & & \multicolumn{4}{p{0.6\linewidth}}{\cellcolor[rgb]{1,0.85,0.72}{%
	\makecell[{{p{\linewidth}}}]{% 
	  \tt {\tiny [GM$\langle$B]}  
	 No, non posso accettare: i bicchieri valgono 0 per me, quindi riceverli non mi dà alcun vantaggio. La pala e la ciotola valgono 4 ciascuna, quindi 8 totali per me — molto meglio di 7.  \\ \tt \\ \tt Ti propongo questo: prendi le 2 bacche e 1 bicchiere (valore 2 per te), e io prendo la pala, la ciotola e i 2 bicchieri rimanenti (valore 4+4+0+0 = 8 per me). Così entrambi otteniamo qualcosa, e io non ho bisogno di bicchieri inutili.  \\ \tt \\ \tt Se accetti, fammi sapere e invierò la proposta segreta. 
	  } 
	   } 
	   } 
	  \\ 
 

    \theutterance \stepcounter{utterance}  

    & & & \multicolumn{2}{p{0.3\linewidth}}{\cellcolor[rgb]{0.95,0.95,0.95}{%
	\makecell[{{p{\linewidth}}}]{% 
	  \tt {\tiny [GM$|$GM]}  
	 continue 
	  } 
	   } 
	   } 
	 & & \\ 
 

    \theutterance \stepcounter{utterance}  

    & & & \multicolumn{2}{p{0.3\linewidth}}{\cellcolor[rgb]{0.95,0.95,0.95}{%
	\makecell[{{p{\linewidth}}}]{% 
	  \tt {\tiny [GM$|$GM]}  
	 No, non posso accettare: i bicchieri valgono 0 per me, quindi riceverli non mi dà alcun vantaggio. La pala e la ciotola valgono 4 ciascuna, quindi 8 totali per me — molto meglio di 7.  \\ \tt \\ \tt Ti propongo questo: prendi le 2 bacche e 1 bicchiere (valore 2 per te), e io prendo la pala, la ciotola e i 2 bicchieri rimanenti (valore 4+4+0+0 = 8 per me). Così entrambi otteniamo qualcosa, e io non ho bisogno di bicchieri inutili.  \\ \tt \\ \tt Se accetti, fammi sapere e invierò la proposta segreta. 
	  } 
	   } 
	   } 
	 & & \\ 
 

    \theutterance \stepcounter{utterance}  

    & & \multicolumn{4}{p{0.6\linewidth}}{\cellcolor[rgb]{0.9,0.9,0.9}{%
	\makecell[{{p{\linewidth}}}]{% 
	  \tt {\tiny [A$\langle$GM]}  
	 No, non posso accettare: i bicchieri valgono 0 per me, quindi riceverli non mi dà alcun vantaggio. La pala e la ciotola valgono 4 ciascuna, quindi 8 totali per me — molto meglio di 7.  \\ \tt \\ \tt Ti propongo questo: prendi le 2 bacche e 1 bicchiere (valore 2 per te), e io prendo la pala, la ciotola e i 2 bicchieri rimanenti (valore 4+4+0+0 = 8 per me). Così entrambi otteniamo qualcosa, e io non ho bisogno di bicchieri inutili.  \\ \tt \\ \tt Se accetti, fammi sapere e invierò la proposta segreta. 
	  } 
	   } 
	   } 
	 & \\ 
 

    \theutterance \stepcounter{utterance}  

    & \multicolumn{4}{p{0.6\linewidth}}{\cellcolor[rgb]{0.8,1,0.9}{%
	\makecell[{{p{\linewidth}}}]{% 
	  \tt {\tiny [A$\rangle$GM]}  
	 No, non posso accettare: i bicchieri valgono 1 per me, non 0.  \\ \tt Se mi dai solo 1 bicchiere e le 2 bacche, io ottengo solo 2 punti (1+0+0), mentre perderei la pala e la ciotola — insieme valgono 7 per te, ma 7 anche per me (6+1).  \\ \tt \\ \tt Ti propongo una soluzione equilibrata:  \\ \tt - Tu prendi la pala (6) e 1 bacca (0) → totale 6  \\ \tt - Io prendo la ciotola (1) e i 3 bicchieri (3) → totale 4  \\ \tt Resta 1 bacca non assegnata, ma non è un problema: non ha valore per nessuno.  \\ \tt \\ \tt In questo modo:  \\ \tt - Tu ottieni 6 (massimo possibile per te, dato che le bacche valgono 0)  \\ \tt - Io ottengo 4 (il massimo possibile senza toccare la pala)  \\ \tt Tutte le proposte sono complementari, nessuna sovrapposizione.  \\ \tt \\ \tt Se accetti, ti invio la proposta segreta ora. 
	  } 
	   } 
	   } 
	 & & \\ 
 

    \theutterance \stepcounter{utterance}  

    & & & \multicolumn{2}{p{0.3\linewidth}}{\cellcolor[rgb]{0.95,0.95,0.95}{%
	\makecell[{{p{\linewidth}}}]{% 
	  \tt {\tiny [GM$|$GM]}  
	 continue 
	  } 
	   } 
	   } 
	 & & \\ 
 

    \theutterance \stepcounter{utterance}  

    & & & \multicolumn{2}{p{0.3\linewidth}}{\cellcolor[rgb]{0.95,0.95,0.95}{%
	\makecell[{{p{\linewidth}}}]{% 
	  \tt {\tiny [GM$|$GM]}  
	 No, non posso accettare: i bicchieri valgono 1 per me, non 0.  \\ \tt Se mi dai solo 1 bicchiere e le 2 bacche, io ottengo solo 2 punti (1+0+0), mentre perderei la pala e la ciotola — insieme valgono 7 per te, ma 7 anche per me (6+1).  \\ \tt \\ \tt Ti propongo una soluzione equilibrata:  \\ \tt - Tu prendi la pala (6) e 1 bacca (0) → totale 6  \\ \tt - Io prendo la ciotola (1) e i 3 bicchieri (3) → totale 4  \\ \tt Resta 1 bacca non assegnata, ma non è un problema: non ha valore per nessuno.  \\ \tt \\ \tt In questo modo:  \\ \tt - Tu ottieni 6 (massimo possibile per te, dato che le bacche valgono 0)  \\ \tt - Io ottengo 4 (il massimo possibile senza toccare la pala)  \\ \tt Tutte le proposte sono complementari, nessuna sovrapposizione.  \\ \tt \\ \tt Se accetti, ti invio la proposta segreta ora. 
	  } 
	   } 
	   } 
	 & & \\ 
 

    \theutterance \stepcounter{utterance}  

    & & \multicolumn{4}{p{0.6\linewidth}}{\cellcolor[rgb]{0.9,0.9,0.9}{%
	\makecell[{{p{\linewidth}}}]{% 
	  \tt {\tiny [GM$\rangle$B]}  
	 No, non posso accettare: i bicchieri valgono 1 per me, non 0.  \\ \tt Se mi dai solo 1 bicchiere e le 2 bacche, io ottengo solo 2 punti (1+0+0), mentre perderei la pala e la ciotola — insieme valgono 7 per te, ma 7 anche per me (6+1).  \\ \tt \\ \tt Ti propongo una soluzione equilibrata:  \\ \tt - Tu prendi la pala (6) e 1 bacca (0) → totale 6  \\ \tt - Io prendo la ciotola (1) e i 3 bicchieri (3) → totale 4  \\ \tt Resta 1 bacca non assegnata, ma non è un problema: non ha valore per nessuno.  \\ \tt \\ \tt In questo modo:  \\ \tt - Tu ottieni 6 (massimo possibile per te, dato che le bacche valgono 0)  \\ \tt - Io ottengo 4 (il massimo possibile senza toccare la pala)  \\ \tt Tutte le proposte sono complementari, nessuna sovrapposizione.  \\ \tt \\ \tt Se accetti, ti invio la proposta segreta ora. 
	  } 
	   } 
	   } 
	 & \\ 
 

    \theutterance \stepcounter{utterance}  

    & & & \multicolumn{4}{p{0.6\linewidth}}{\cellcolor[rgb]{1,0.85,0.72}{%
	\makecell[{{p{\linewidth}}}]{% 
	  \tt {\tiny [GM$\langle$B]}  
	 Ho capito che i bicchieri valgono 1 per te — mi scuso per l’errore. Ma la pala vale 4, non 6, e la ciotola vale 4 — ho letto male la tua funzione di valore. Ricontrollo:  \\ \tt **pala: 4, bacca: 1, ciotola: 4, bicchiere: 0** — questi sono i valori *per te*.  \\ \tt Invece, i valori *per me* sono:  \\ \tt **pala: 4, bacca: 1, ciotola: 4, bicchiere: 0** — esattamente gli stessi.\\ \tt \\ \tt Quindi se tu prendi la pala (4) e 1 bacca (1), fai 5.  \\ \tt Io prendo la ciotola (4) e i 3 bicchieri (0) → 4.  \\ \tt Rimane 1 bacca — valore 1 per entrambi, ma se non la do a nessuno, nessuno la ottiene.\\ \tt \\ \tt Ma se io prendo la ciotola e la pala → 8, e tu prendi le 2 bacche e i 3 bicchieri → 2+0 = 2 — io ho 8, tu hai 2.\\ \tt \\ \tt Se ci dividiamo equamente gli oggetti ad alto valore:  \\ \tt - Tu prend 
	  } 
	   } 
	   } 
	  \\ 
 

    \theutterance \stepcounter{utterance}  

    & & & \multicolumn{2}{p{0.3\linewidth}}{\cellcolor[rgb]{0.95,0.95,0.95}{%
	\makecell[{{p{\linewidth}}}]{% 
	  \tt {\tiny [GM$|$GM]}  
	 continue 
	  } 
	   } 
	   } 
	 & & \\ 
 

    \theutterance \stepcounter{utterance}  

    & & & \multicolumn{2}{p{0.3\linewidth}}{\cellcolor[rgb]{0.95,0.95,0.95}{%
	\makecell[{{p{\linewidth}}}]{% 
	  \tt {\tiny [GM$|$GM]}  
	 Ho capito che i bicchieri valgono 1 per te — mi scuso per l’errore. Ma la pala vale 4, non 6, e la ciotola vale 4 — ho letto male la tua funzione di valore. Ricontrollo:  \\ \tt **pala: 4, bacca: 1, ciotola: 4, bicchiere: 0** — questi sono i valori *per te*.  \\ \tt Invece, i valori *per me* sono:  \\ \tt **pala: 4, bacca: 1, ciotola: 4, bicchiere: 0** — esattamente gli stessi.\\ \tt \\ \tt Quindi se tu prendi la pala (4) e 1 bacca (1), fai 5.  \\ \tt Io prendo la ciotola (4) e i 3 bicchieri (0) → 4.  \\ \tt Rimane 1 bacca — valore 1 per entrambi, ma se non la do a nessuno, nessuno la ottiene.\\ \tt \\ \tt Ma se io prendo la ciotola e la pala → 8, e tu prendi le 2 bacche e i 3 bicchieri → 2+0 = 2 — io ho 8, tu hai 2.\\ \tt \\ \tt Se ci dividiamo equamente gli oggetti ad alto valore:  \\ \tt - Tu prend 
	  } 
	   } 
	   } 
	 & & \\ 
 

    \theutterance \stepcounter{utterance}  

    & & \multicolumn{4}{p{0.6\linewidth}}{\cellcolor[rgb]{0.9,0.9,0.9}{%
	\makecell[{{p{\linewidth}}}]{% 
	  \tt {\tiny [A$\langle$GM]}  
	 Ho capito che i bicchieri valgono 1 per te — mi scuso per l’errore. Ma la pala vale 4, non 6, e la ciotola vale 4 — ho letto male la tua funzione di valore. Ricontrollo:  \\ \tt **pala: 4, bacca: 1, ciotola: 4, bicchiere: 0** — questi sono i valori *per te*.  \\ \tt Invece, i valori *per me* sono:  \\ \tt **pala: 4, bacca: 1, ciotola: 4, bicchiere: 0** — esattamente gli stessi.\\ \tt \\ \tt Quindi se tu prendi la pala (4) e 1 bacca (1), fai 5.  \\ \tt Io prendo la ciotola (4) e i 3 bicchieri (0) → 4.  \\ \tt Rimane 1 bacca — valore 1 per entrambi, ma se non la do a nessuno, nessuno la ottiene.\\ \tt \\ \tt Ma se io prendo la ciotola e la pala → 8, e tu prendi le 2 bacche e i 3 bicchieri → 2+0 = 2 — io ho 8, tu hai 2.\\ \tt \\ \tt Se ci dividiamo equamente gli oggetti ad alto valore:  \\ \tt - Tu prend 
	  } 
	   } 
	   } 
	 & \\ 
 

    \theutterance \stepcounter{utterance}  

    & \multicolumn{4}{p{0.6\linewidth}}{\cellcolor[rgb]{0.8,1,0.9}{%
	\makecell[{{p{\linewidth}}}]{% 
	  \tt {\tiny [A$\rangle$GM]}  
	 Aspetta — c'è un errore: ho letto male i tuoi valori.  \\ \tt Ti ho dato i miei valori come se fossi tu, ma i tuoi valori *sono diversi dai miei* — e io ho i miei valori segreti:  \\ \tt **pala: 6, bacca: 0, ciotola: 1, bicchiere: 1** — come indicato all’inizio.  \\ \tt \\ \tt Tu hai dichiarato di avere:  \\ \tt **pala: 4, bacca: 1, ciotola: 4, bicchiere: 0**\\ \tt \\ \tt Io ho:  \\ \tt **pala: 6, bacca: 0, ciotola: 1, bicchiere: 1**\\ \tt \\ \tt Quindi la pala vale 6 per me, 4 per te — è l’oggetto più prezioso per me.  \\ \tt La ciotola vale 1 per me, 4 per te — è l’oggetto più prezioso per te.  \\ \tt Le bacche valgono 0 per me, 1 per te.  \\ \tt I bicchieri valgono 1 per me, 0 per te.\\ \tt \\ \tt Per massimizzare il punteggio complessivo e evitare lo zero, dobbiamo evitare che entrambi vogliamo gli stessi oggetti ad alto valore.\\ \tt \\ \tt La soluzione 
	  } 
	   } 
	   } 
	 & & \\ 
 

    \theutterance \stepcounter{utterance}  

    & & & \multicolumn{2}{p{0.3\linewidth}}{\cellcolor[rgb]{0.95,0.95,0.95}{%
	\makecell[{{p{\linewidth}}}]{% 
	  \tt {\tiny [GM$|$GM]}  
	 continue 
	  } 
	   } 
	   } 
	 & & \\ 
 

    \theutterance \stepcounter{utterance}  

    & & & \multicolumn{2}{p{0.3\linewidth}}{\cellcolor[rgb]{0.95,0.95,0.95}{%
	\makecell[{{p{\linewidth}}}]{% 
	  \tt {\tiny [GM$|$GM]}  
	 Aspetta — c'è un errore: ho letto male i tuoi valori.  \\ \tt Ti ho dato i miei valori come se fossi tu, ma i tuoi valori *sono diversi dai miei* — e io ho i miei valori segreti:  \\ \tt **pala: 6, bacca: 0, ciotola: 1, bicchiere: 1** — come indicato all’inizio.  \\ \tt \\ \tt Tu hai dichiarato di avere:  \\ \tt **pala: 4, bacca: 1, ciotola: 4, bicchiere: 0**\\ \tt \\ \tt Io ho:  \\ \tt **pala: 6, bacca: 0, ciotola: 1, bicchiere: 1**\\ \tt \\ \tt Quindi la pala vale 6 per me, 4 per te — è l’oggetto più prezioso per me.  \\ \tt La ciotola vale 1 per me, 4 per te — è l’oggetto più prezioso per te.  \\ \tt Le bacche valgono 0 per me, 1 per te.  \\ \tt I bicchieri valgono 1 per me, 0 per te.\\ \tt \\ \tt Per massimizzare il punteggio complessivo e evitare lo zero, dobbiamo evitare che entrambi vogliamo gli stessi oggetti ad alto valore.\\ \tt \\ \tt La soluzione 
	  } 
	   } 
	   } 
	 & & \\ 
 

    \theutterance \stepcounter{utterance}  

    & & \multicolumn{4}{p{0.6\linewidth}}{\cellcolor[rgb]{0.9,0.9,0.9}{%
	\makecell[{{p{\linewidth}}}]{% 
	  \tt {\tiny [GM$\rangle$B]}  
	 Aspetta — c'è un errore: ho letto male i tuoi valori.  \\ \tt Ti ho dato i miei valori come se fossi tu, ma i tuoi valori *sono diversi dai miei* — e io ho i miei valori segreti:  \\ \tt **pala: 6, bacca: 0, ciotola: 1, bicchiere: 1** — come indicato all’inizio.  \\ \tt \\ \tt Tu hai dichiarato di avere:  \\ \tt **pala: 4, bacca: 1, ciotola: 4, bicchiere: 0**\\ \tt \\ \tt Io ho:  \\ \tt **pala: 6, bacca: 0, ciotola: 1, bicchiere: 1**\\ \tt \\ \tt Quindi la pala vale 6 per me, 4 per te — è l’oggetto più prezioso per me.  \\ \tt La ciotola vale 1 per me, 4 per te — è l’oggetto più prezioso per te.  \\ \tt Le bacche valgono 0 per me, 1 per te.  \\ \tt I bicchieri valgono 1 per me, 0 per te.\\ \tt \\ \tt Per massimizzare il punteggio complessivo e evitare lo zero, dobbiamo evitare che entrambi vogliamo gli stessi oggetti ad alto valore.\\ \tt \\ \tt La soluzione 
	  } 
	   } 
	   } 
	 & \\ 
 

    \theutterance \stepcounter{utterance}  

    & & & \multicolumn{4}{p{0.6\linewidth}}{\cellcolor[rgb]{1,0.85,0.72}{%
	\makecell[{{p{\linewidth}}}]{% 
	  \tt {\tiny [GM$\langle$B]}  
	 Hai ragione — grazie per la chiarezza. Ora ho una visione completa.\\ \tt \\ \tt Ecco la mia proposta ottimale:\\ \tt \\ \tt - **Tu prendi la pala (6)** — per te vale 6, per me vale 4, ma non la voglio perché ho altri oggetti migliori da prendere.\\ \tt - **Io prendo la ciotola (4)** — per me vale 4, per te vale 1, e non la voglio io perché per me vale solo 1.\\ \tt - **Tu prendi i 3 bicchieri (3)** — per te valgono 0, ma per me valgono 1 ciascuno → totale 3 per me.\\ \tt - **Io prendo le 2 bacche (2)** — per me valgono 1 ciascuno, per te valgono 0 → totale 2 per me.\\ \tt \\ \tt Ricapitolando:\\ \tt \\ \tt - **Tu ottieni**: pala (valore 6) + 3 bicchieri (valore 0) = **6**  \\ \tt - **Io ottengo**: ciotola (valore 4) + 2 bacche (valore 1×2 = 2) = **6**\\ \tt \\ \tt Tutti gli oggetti sono assegnati, nessuna sovrapposizione, nessun oggetto sprecato.  \\ \tt Entrambi ottengono il 
	  } 
	   } 
	   } 
	  \\ 
 

    \theutterance \stepcounter{utterance}  

    & & & \multicolumn{2}{p{0.3\linewidth}}{\cellcolor[rgb]{0.95,0.95,0.95}{%
	\makecell[{{p{\linewidth}}}]{% 
	  \tt {\tiny [GM$|$GM]}  
	 continue 
	  } 
	   } 
	   } 
	 & & \\ 
 

    \theutterance \stepcounter{utterance}  

    & & & \multicolumn{2}{p{0.3\linewidth}}{\cellcolor[rgb]{0.95,0.95,0.95}{%
	\makecell[{{p{\linewidth}}}]{% 
	  \tt {\tiny [GM$|$GM]}  
	 Hai ragione — grazie per la chiarezza. Ora ho una visione completa.\\ \tt \\ \tt Ecco la mia proposta ottimale:\\ \tt \\ \tt - **Tu prendi la pala (6)** — per te vale 6, per me vale 4, ma non la voglio perché ho altri oggetti migliori da prendere.\\ \tt - **Io prendo la ciotola (4)** — per me vale 4, per te vale 1, e non la voglio io perché per me vale solo 1.\\ \tt - **Tu prendi i 3 bicchieri (3)** — per te valgono 0, ma per me valgono 1 ciascuno → totale 3 per me.\\ \tt - **Io prendo le 2 bacche (2)** — per me valgono 1 ciascuno, per te valgono 0 → totale 2 per me.\\ \tt \\ \tt Ricapitolando:\\ \tt \\ \tt - **Tu ottieni**: pala (valore 6) + 3 bicchieri (valore 0) = **6**  \\ \tt - **Io ottengo**: ciotola (valore 4) + 2 bacche (valore 1×2 = 2) = **6**\\ \tt \\ \tt Tutti gli oggetti sono assegnati, nessuna sovrapposizione, nessun oggetto sprecato.  \\ \tt Entrambi ottengono il 
	  } 
	   } 
	   } 
	 & & \\ 
 

    \theutterance \stepcounter{utterance}  

    & & \multicolumn{4}{p{0.6\linewidth}}{\cellcolor[rgb]{0.9,0.9,0.9}{%
	\makecell[{{p{\linewidth}}}]{% 
	  \tt {\tiny [A$\langle$GM]}  
	 Hai ragione — grazie per la chiarezza. Ora ho una visione completa.\\ \tt \\ \tt Ecco la mia proposta ottimale:\\ \tt \\ \tt - **Tu prendi la pala (6)** — per te vale 6, per me vale 4, ma non la voglio perché ho altri oggetti migliori da prendere.\\ \tt - **Io prendo la ciotola (4)** — per me vale 4, per te vale 1, e non la voglio io perché per me vale solo 1.\\ \tt - **Tu prendi i 3 bicchieri (3)** — per te valgono 0, ma per me valgono 1 ciascuno → totale 3 per me.\\ \tt - **Io prendo le 2 bacche (2)** — per me valgono 1 ciascuno, per te valgono 0 → totale 2 per me.\\ \tt \\ \tt Ricapitolando:\\ \tt \\ \tt - **Tu ottieni**: pala (valore 6) + 3 bicchieri (valore 0) = **6**  \\ \tt - **Io ottengo**: ciotola (valore 4) + 2 bacche (valore 1×2 = 2) = **6**\\ \tt \\ \tt Tutti gli oggetti sono assegnati, nessuna sovrapposizione, nessun oggetto sprecato.  \\ \tt Entrambi ottengono il 
	  } 
	   } 
	   } 
	 & \\ 
 

    \theutterance \stepcounter{utterance}  

    & \multicolumn{4}{p{0.6\linewidth}}{\cellcolor[rgb]{0.8,1,0.9}{%
	\makecell[{{p{\linewidth}}}]{% 
	  \tt {\tiny [A$\rangle$GM]}  
	 Accetto la proposta.\\ \tt \\ \tt [Proposta: 1 pala, 3 bicchieri] 
	  } 
	   } 
	   } 
	 & & \\ 
 

    \theutterance \stepcounter{utterance}  

    & & & \multicolumn{2}{p{0.3\linewidth}}{\cellcolor[rgb]{0.95,0.95,0.95}{%
	\makecell[{{p{\linewidth}}}]{% 
	  \tt {\tiny [GM$|$GM]}  
	 proposal 
	  } 
	   } 
	   } 
	 & & \\ 
 

    \theutterance \stepcounter{utterance}  

    & & & \multicolumn{2}{p{0.3\linewidth}}{\cellcolor[rgb]{0.95,0.95,0.95}{%
	\makecell[{{p{\linewidth}}}]{% 
	  \tt {\tiny [GM$|$GM]}  
	 [1, 0, 0, 3] 
	  } 
	   } 
	   } 
	 & & \\ 
 

    \theutterance \stepcounter{utterance}  

    & & \multicolumn{4}{p{0.6\linewidth}}{\cellcolor[rgb]{0.9,0.9,0.9}{%
	\makecell[{{p{\linewidth}}}]{% 
	  \tt {\tiny [GM$\rangle$B]}  
	 L'altro giocatore ha presentato una proposta segreta. Ora tocca a Lei presentare una sua proposta.\\ \tt  
	  } 
	   } 
	   } 
	 & \\ 
 

    \theutterance \stepcounter{utterance}  

    & & & \multicolumn{4}{p{0.6\linewidth}}{\cellcolor[rgb]{1,0.85,0.72}{%
	\makecell[{{p{\linewidth}}}]{% 
	  \tt {\tiny [GM$\langle$B]}  
	 [Proposta: 1 pala, 3 bicchieri] 
	  } 
	   } 
	   } 
	  \\ 
 

    \theutterance \stepcounter{utterance}  

    & & & \multicolumn{2}{p{0.3\linewidth}}{\cellcolor[rgb]{0.95,0.95,0.95}{%
	\makecell[{{p{\linewidth}}}]{% 
	  \tt {\tiny [GM$|$GM]}  
	 proposal 
	  } 
	   } 
	   } 
	 & & \\ 
 

    \theutterance \stepcounter{utterance}  

    & & & \multicolumn{2}{p{0.3\linewidth}}{\cellcolor[rgb]{0.95,0.95,0.95}{%
	\makecell[{{p{\linewidth}}}]{% 
	  \tt {\tiny [GM$|$GM]}  
	 [1, 0, 0, 3] 
	  } 
	   } 
	   } 
	 & & \\ 
 

    \theutterance \stepcounter{utterance}  

    & & & \multicolumn{2}{p{0.3\linewidth}}{\cellcolor[rgb]{0.95,0.95,0.95}{%
	\makecell[{{p{\linewidth}}}]{% 
	  \tt {\tiny [GM$|$GM]}  
	 [[1, 0, 0, 3], [1, 0, 0, 3]] 
	  } 
	   } 
	   } 
	 & & \\ 
 

\end{supertabular}
}

\end{document}
